\section{Investigación}

El manejo de los datos en un proyecto, trabajo, artículo, investigación, etc. es una parte fundamental ya que vivimos en una era en la que cualquier dato que pueda brindar alguna mínima información sobre alguien puede llegar a ser utilizado en su contra, perjudicandolo cuando esa persona no realizó la investigación y no formó parte de su formulación, aún así se vió afectada. Es por ello que los encargados de los datos tienen que asegurarse de que tienen el permiso de usarlos y en caso de que muestren datos sensibles, anonimizarlos o simplemente no compartirlos para no involucrar a terceros.

Teniendo en cuenta esto es que de igual manera en este proyecto es necesario tomar precauciones sobre los datos a pesar de que se tratan de vacas, siempre puede llegar a presentarse un conflicto con su publicación, por ello es que se realizó una investigación sobre qué normas o estándares se utilizan en la comunidad científica para seguir sus mismos pasos y no presentar problemas.

\subsection{ISO/IEC 25012:2008}

La primera norma que seguimos es esta, es parte de la familia SQuaRE. Define un modelo general de calidad para datos guardados de forma estructurada en un sistema computacional y sirve para establecer requisitos, definir medidas y planear evaluaciones de calidad de datos. Entre sus usos, menciona evaluar si los datos cumplen con la legislación o con ciertos requisitos.

Agrupa los atributos de calidad en quince características, vistas desde dos puntos de vista: inherentes y dependientes del sistema.

\begin{itemize}
  \item Inherentes: exactitud, completitud, consistencia, credibilidad, actualidad.
  \item Inherentes y del sistema: accesibilidad, cumplimiento, confidencialidad, eficiencia, precisión, trazabilidad, comprensibilidad.
  \item Del sistema: disponibilidad, portabilidad, recuperabilidad.
\end{itemize}

La manera en que esta norma se puede aplicar en el proyecto es por medio de la producción diaria de leche, ya que revisa la completitud (días sin datos), consistencia (unidades y formatos) y actualidad. También en las fotos y etiquetas de BCS porque podemos medir la exactitud y precisión (pasos de 0.25), credibilidad (quién calificó), consistencia (ángulo, altura, distancia) y trazabilidad (que cada foto quede ligada a Cow ID y fecha).

\subsection{ISO/IEC 29100:2024}

Por otro lado, tenemos otra norma pero en este caso sobre la privacidad en el uso de los datos, define los actores y sus roles al procesar información personal identificable (PII), describe consideraciones de protección y remite a principios de privacidad conocidos. En la página de ISO aparece con costo de CHF 0, así que puede consultarse gratis.

Define roles (la persona titular de los datos, quien decide el tratamiento y quien lo ejecuta) en la primera edición se cuentan con 11 principios: consentimiento y elección, legitimidad y especificación del propósito, limitación de la recolección, minimización de datos, limitación de uso, retención y divulgación, exactitud y calidad, transparencia, participación individual, responsabilidad, seguridad de la información y cumplimiento.

Esta norma se centra en datos personales, por lo que podría pensarse en primera instancia que a las vacas no podría aplicar, pero la realidad es otra ya que puede ponerse en práctica debido a nombres de empleados, ordeñadores o propietarios en las tablas; además de personas que salgan en segundo plano en las fotos.

Debido a esta norma es que es importante considerar seguir una serie de principios para poder tener un orden con el manejo de los datos, además de evitar tener alguna información personal fuera de lugar en donde cualquiera pueda hacer mal uso de ella.

\begin{itemize}
  \item Propósito: Usar los datos solamente para el proyecto académico y evitar usarlos en otros contextos no seguros.
  \item Minimización: Quedarse únicamente con las columnas necesarias.
  \item Retención: Borrar o devolver los datos al terminar, si es que el rancho lo solicita.
  \item Responsabilidad: Nombrar quién del equipo custodiará los datos.
  \item Seguridad: No subirlos a repositorios públicos ni servicios externos sin permiso.
\end{itemize}

Algunos papers en donde podemos encontrar el seguimiento de estas normas son los siguientes:

\begin{itemize}
  \item H. A. J. Candelario, T. V. Cruz, P. A. R. Gray, A. A. R. C. Sison, C. J. Centeno and J. S. Diego, “Web-Based Application Predicting Cow Hoof Infections with 3D Modeling Scrapping Simulator and Treatment Recommendation,” 2024 International Conference on Intelligent Cybernetics Technology \& Applications (ICICyTA), Bali, Indonesia, 2024, pp. 937-942, doi: 10.1109/ICICYTA64807.2024.10913372.
  \item Bewong, M., Ip, R. H. L., Lewis, C., Krivokapic-Skoko, B., Islam, Z., Al-Saggaf, Y., Medway, J., Ali, B., \& Dixon, E. (2023). Potential implications and benefits for the agricultural technology sector from the introduction of the Australian Agricultural Data Exchange. Food Agility Cooperative Centre.
\end{itemize}

Gracias a estos artículos y las normas que siguen, es que podemos darnos una idea de lo que es necesario o deberíamos de implementar en el proyecto con los datos proporcionados y adquiridos. Muchos o la gran mayoría de los artículos dedicados a la ganadería no presentan ninguna declaración sobre el uso de los datos, simplemente mencionan agradecimientos al rancho del cual los obtuvieron pero nada más, por lo que este punto también es importante para nosotros, agradecer al rancho CAETEC que nos permitió hacer uso de sus datos, pero nosotros no lo dejaremos simplemente en este punto, sino también complementar con lo que encontramos en los papers mencionados anteriormente.

Cumplir con estas normas se ha vuelto parte esencial de un trabajo en el que se involucran datos, por ello al no encontrar en primera instancia papers que mencionan su uso de datos, la forma en la que se trabajó fue de primera instancia buscar normas generales sobre protecciones de datos y ya una vez seleccionadas las que nos parecieron más importantes y acordes con el proyecto, es que se buscaron los papers teniendo en cuenta que hablen sobre su uso de datos siguiendo estas normas, y sobre todo que sean de un área relacionada, ya sea que trabajen explícitamente con vacas o simplemente que sea un trabajo relacionado a la ganadería.
